% Kapitel 1: Einleitung
\section{Einleitung}

Das Immersive Reality Lab (IRL) der Hochschule Hamm-Lippstadt bündelt auf \url{https://immersive-reality-lab.de} Informationen zu Forschung, Projekten, Publikationen und Lehre. Mit wachsendem Umfang steigt der Aufwand, gezielt Informationen zu finden.

\textbf{Human-Centred Design (HCD)} nach DIN EN ISO 9241-210~\cite{iso9241-210} ist das zentrale Vorgehensmodell dieses Projekts: Nutzer werden früh einbezogen, Anforderungen werden aus Nutzungssituation und User Research abgeleitet, der Prototyp wird evaluiert und iterativ verbessert. Technische Entscheidungen (RAG, Chunking, Fehlertexte) werden ebenso gegen ergonomische Kriterien geprüft wie die Oberfläche.

Ziel von \textit{Ergonomia} ist ein RAG-basierter Chatbot für das IRL. Die Bewertung stützt sich auf Amershi et al.~\cite{amershi2019}, Nielsen~\cite{nielsen1994}, DIN EN ISO 9241-110~\cite{iso9241} und den SUS~\cite{lewis2018}.
